% 行星（综述）
% license CCBYSA3
% type Wiki

本文根据 CC-BY-SA 协议转载翻译自维基百科\href{https://en.wikipedia.org/wiki/Planet}{相关文章}。

\begin{figure}[ht]
\centering
\includegraphics[width=8cm]{./figures/c14ce525d8184ebd.png}
\caption{太阳系八大行星的尺寸比例图（自上而下、从左到右）：土星、木星、天王星、海王星（外行星），地球、金星、火星和水星（内行星）。} \label{fig_Planet_1}
\end{figure}
行星是一种巨大且呈圆形的天体，一般要求其绕恒星、恒星遗迹或褐矮星运行，并且其自身不属于这些天体。按该术语最严格的定义，太阳系共有八颗行星：类地行星（水星、金星、地球和火星）以及巨行星（木星、土星、天王星和海王星）。关于行星形成的最佳理论是星云假说，该假说认为一团星际云从星云中坍缩，形成一颗年轻的原恒星，而这颗原恒星被一个原行星盘所环绕。行星在该盘内通过由引力驱动的物质逐渐累积而成长，这一过程称为吸积。

“行星”（planet）一词源于希腊语 πλανήται（planḗtai），意为“流浪者”。在古代，这个词指的是太阳、月亮，以及在恒星背景上肉眼可见并会移动的五个光点——水星、金星、火星、木星和土星。行星在历史上具有宗教关联：多种文化将天体与神祇对应，而这种与神话和民间传说的联系，也延续到了对新发现的太阳系天体的命名方式之中。随着16至17世纪日心说取代地心说，地球本身也被认作一颗行星。

随着望远镜的发展，“行星”一词的含义被扩展，纳入了只能借助观测工具才能看见的天体：地球之外行星的卫星；冰巨行星天王星与海王星；谷神星及后来被确认属于小行星带的其他天体；以及冥王星，后来被发现是柯伊伯带中冰质天体族群中最大的成员。对柯伊伯带中其他大型天体的发现，尤其是厄里斯，引发了关于如何精确定义“行星”的争论。2006年，国际天文学联合会（IAU）通过了太阳系中“行星”的定义，将四颗类地行星与四颗巨行星归入行星范畴；而谷神星、冥王星与厄里斯则被归入“矮行星”类别。尽管如此，许多行星科学家仍持续以更宽泛的方式使用“行星”一词，包括矮行星以及如月球般呈圆形的卫星。

天文学的进一步发展使得人们发现了超过5900颗太阳系外的行星，即系外行星。这些系外行星常呈现太阳系行星所不具备的特殊性质，例如“热木星”——如飞马座51b般绕其母恒星极近距离运行的巨行星——以及极端偏心的轨道，如HD~20782~b。褐矮星与比木星更大的行星的发现也推动了对于定义边界的讨论，即如何划定行星与恒星之间的分界线。已有多颗系外行星被发现运行在其恒星的宜居带内（即行星表面可能存在液态水的区域），但地球仍是唯一已知能够承载生命的行星。宜居带的形成受多种因素影响，例如行星的大气组成（若其具有大气）。例如火星在赤道区域的白天温度足以形成液态水，只是其气压过低而无法实现。
\subsection{形成}
\begin{figure}[ht]
\centering
\includegraphics[width=8cm]{./figures/edc16278f48809b3.png}
\caption{原行星盘} \label{fig_Planet_2}
\end{figure}
\begin{figure}[ht]
\centering
\includegraphics[width=8cm]{./figures/645312c0cec1ae7e.png}
\caption{行星形成过程中原行星之间的碰撞} \label{fig_Planet_4}
\end{figure}
行星的形成机制目前尚未被完全确定。主流理论认为，行星在星云塌缩形成一层薄的气体和尘埃盘的过程中逐渐聚合而成。在中心区域形成原恒星，其周围则是一个旋转的原行星盘。通过吸积（即通过“粘性碰撞”的过程），盘中的尘埃粒子逐渐聚集质量，形成越来越大的天体。在局部区域中会出现称为微行星的质量集中体，它们通过自身的引力吸引更多物质，从而加速吸积过程。这些质量集中体变得越来越致密，最终在自引力作用下向内塌缩，形成原行星 [6]。

当一个原行星的质量增长到略高于火星的质量时，它会开始吸积一层延展的大气层 [7]，从而通过大气阻力效应 [8][9] 极大提高其对微行星的捕获效率。根据固体与气体的吸积历史不同，最终可能形成巨行星、冰巨行星或类地行星 [10][11][12]。人们认为木星、土星和天王星的规则卫星可能以类似方式形成 [13][14]；然而，海王星的卫星海卫一很可能是被捕获的 [15]。地球的月球 [16] 与冥王星的卫星卡戎 [17] 则可能是在巨型碰撞中形成。

当原恒星增大至能够点燃并成为真正的恒星时，残余的吸积盘会从内向外被光致蒸发、太阳风、Poynting–Robertson 阻力及其他效应逐渐清除 [18][19]。此后，系统中可能仍存在许多原行星，它们相互环绕或环绕恒星运行，但随着时间推移，其中许多会发生碰撞，或合并成更大的原行星，或为其他原行星提供可吸积的碎片 [20]。那些质量足够大的天体会吸尽其轨道邻域的大部分物质而成为行星。避免了碰撞的原行星可能通过引力捕获成为行星的天然卫星，或者留存在某些轨道带中，最终成为矮行星或小天体 [21][22]。
\begin{figure}[ht]
\centering
\includegraphics[width=8cm]{./figures/624858ff17708327.png}
\caption{超新星遗迹喷 ejecta 产生行星形成物质} \label{fig_Planet_5}
\end{figure}
较小的微行星的高能撞击（以及放射性衰变）会加热正在成长的行星，使其至少部分熔化。行星内部开始按照密度分异，密度更高的物质下沉至核心。[23] 较小的类地行星在这种聚积过程中会失去其大部分大气，但失去的气体可通过地幔逸气以及随后彗星的撞击所补充[24]（更小的行星会通过各种逃逸机制而失去它们获得的任何大气[25]）。

随着对太阳以外恒星的行星系统的发现与观测，人们得以对上述理论加以扩展、修正甚至替代。金属丰度——一个天文学术语，用来描述原子序数大于 2（即氦）的化学元素的丰度——似乎决定着恒星拥有行星的可能性。[26][27] 因此，与金属贫乏的第二星族（Population II）恒星相比，金属丰富的第一星族（Population I）恒星更可能拥有规模可观的行星系统。[28]
\subsection{太阳系中的行星}
根据 IAU 的定义，太阳系中共有八颗行星（按距太阳距离递增排列）[2]：Mercury、Venus、Earth、Mars、Jupiter、Saturn、Uranus 与 Neptune。Jupiter 质量最大，约为地球质量的 318 倍；而 Mercury 最小，仅为地球质量的 0.055 倍。[29]

太阳系行星可依据组成成分划分为不同类别。类地行星（terrestrials）与 Earth 类似，其主体由岩石与金属构成：Mercury、Venus、Earth 与 Mars。Earth 是类地行星中体积最大的一个。[30]

巨行星（giant planets）的质量显著大于类地行星：Jupiter、Saturn、Uranus 与 Neptune。[30] 它们在成分上与类地行星存在根本差异。气体巨行星（gas giants）Jupiter 与 Saturn 主要由氢与氦组成，是太阳系中最为巨大的行星。Saturn 的质量约为 Jupiter 的三分之一，即约 95 个地球质量。[31]

冰巨行星（ice giants）Uranus 与 Neptune 则主要由低沸点物质组成，如 water、methane 与 ammonia，并拥有由 hydrogen 与 helium 构成的厚大气层。它们的质量显著低于气体巨行星，仅分别为 14 与 17 个地球质量。[31]
\begin{figure}[ht]
\centering
\includegraphics[width=8cm]{./figures/f984108e351cc1ba.png}
\caption{太阳、行星、矮行星与卫星的尺寸按比例标注示意。各天体之间的距离不按比例绘制。小行星带位于火星与木星轨道之间，柯伊伯带位于海王星轨道之外。} \label{fig_Planet_6}
\end{figure}
\noindent 矮行星是达到引力塑形（即具有趋于球形的外形），但尚未清除其轨道附近其他天体的天体。按照距太阳平均距离递增排列，被天文学界普遍认可的矮行星包括：Ceres，Orcus，Pluto，Haumea，Quaoar，Makemake，Gonggong，Eris 与 Sedna。[32][33]Ceres 位于 Mars 与 Jupiter 轨道之间的小行星带，是该区域最大的天体。其余八个均位于 Neptune 轨道之外。Orcus、Pluto、Haumea、Quaoar 与 Makemake 运行于 Kuiper 带——这是位于 Neptune 轨道之外的第二个太阳系小天体带。Gonggong 与 Eris 运行于散射盘（scattered disc），其位置更为外侧，且该区域不同于 Kuiper 带——会因 Neptune 的引力扰动而呈现轨道不稳定性。Sedna 是目前已知最大的游离天体（detached object），其轨道不会接近太阳到足以受任何经典行星的显著影响；其轨道成因仍在讨论之中。以上九个天体均与类地行星相似，拥有固体表面，但其成分主要是冰与岩石，而非岩石与金属。此外，它们全部小于 Mercury，其中 Pluto 是目前已知体积最大的矮行星，而 Eris 的质量最大。[34][35]

至少有十九颗具行星质量的卫星（planetary-mass moons 或 satellite planets），即质量足以呈现椭球形状的卫星[4]：

\begin{itemize}
    \item Earth 的一颗卫星：Moon；
    \item Jupiter 的四颗卫星：Io，Europa，Ganymede 与 Callisto；
    \item Saturn 的七颗卫星：Mimas，Enceladus，Tethys，Dione，Rhea，Titan 与 Iapetus；
    \item Uranus 的五颗卫星：Miranda，Ariel，Umbriel，Titania 与 Oberon；
    \item Neptune 的一颗卫星：Triton；
    \item Pluto 的一颗卫星：Charon。
\end{itemize}
Moon、Io 与 Europa 的成分与类地行星相似；其它卫星的成分则与矮行星类似，以冰与岩石为主，其中 Tethys 几乎完全由冰组成。由于其表面的冰层使内部难以观测，Europa 常被视为一颗冰质行星。[4][36]

Ganymede 与 Titan 在半径上均大于 Mercury，而 Callisto 也几乎与之相当，但这三者的质量均远小于 Mercury。Mimas 是公认的几何意义上的行星（即达到地质行星标准）中最小的一个，其质量约为地球质量的六百万分之一；然而，还有许多比它更大的天体却可能并非地质行星（例如 Salacia）。[32]
\subsection{Exoplanets}
\begin{figure}[ht]
\centering
\includegraphics[width=8cm]{./figures/85a4bfec6f369e50.png}
\caption{截至 2023 年 8 月的每年系外行星探测数量（数据来自 NASA 系外行星档案）[37]} \label{fig_Planet_8}
\end{figure}
系外行星（exoplanet）指位于太阳系之外的行星。截至 2025 年 10 月 30 日，已经确认的系外行星共有 6{,}128 颗，分布在 4{,}584 个行星系统中，其中 1{,}017 个系统拥有不止一颗行星。[38]已知的系外行星尺寸跨度极大，从比 Jupiter 大两倍的气态巨行星到仅略大于 Moon 的微小行星均有分布。对引力微透镜（gravitational microlensing）数据的分析表明，在 Milky Way 中，每颗恒星平均至少拥有 1.6 颗束缚行星。[39]

1992 年初，无线电天文学家 Aleksander Wolszczan 与 Dale Frail 宣布发现两颗绕脉冲星 PSR 1257+12 运行的行星。[40] 这项发现随后得到证实，并普遍被认为是首次确定无疑的系外行星探测。研究者推测，这些行星可能形成于产生该脉冲星的超新星爆炸后残留的盘状物质中。[41]

首次确认的、绕普通主序星运行的系外行星于 1995 年 10 月 6 日被发现。当时 Geneva 大学的 Michel Mayor 与 Didier Queloz 宣布在恒星 51 Pegasi 周围探测到行星 51 Pegasi b。[42]从那时直到 Kepler 空间望远镜任务开始前，人类已知的大多数系外行星都是质量与 Jupiter 相当或更大的气态巨行星，因为它们更易被探测到。Kepler 候选行星目录中大多数行星的尺度介于 Neptune 与更小的行星之间，有些甚至小于 Mercury。[43][44]

2011 年，Kepler 团队报告了首次发现绕类太阳恒星运行的地球大小行星：Kepler-20e 与 Kepler-20f。[45][46][47]自那时起，已有超过 100 颗大小近似 Earth 的行星被确认，其中有 20 颗位于其恒星的宜居带（habitable zone）内——这是允许类地行星在具备足够大气压的条件下，其表面可能存在液态水的轨道范围。[48][49][50]据估计，每五颗类太阳恒星中就有一颗拥有位于宜居带中的地球大小行星，这意味着距离地球最近的此类行星预计不会超过 12 光年。[a]此类类地行星的出现频率是 Drake 方程中的变量之一，该方程用于估算 Milky Way 中存在的具备智慧并能够通信的文明数量。[53]

在太阳系中不存在的行星类型包括超地球（super-Earth）与迷你海王星（mini-Neptune），其质量介于 Earth 与 Neptune 之间。质量不超过约两倍地球质量的天体通常被认为具有类似 Earth 的岩石组成；超过该范围，则会演化为含有大量挥发物与气体、类似 Neptune 的结构。[54]行星 Gliese~581c 的质量为 Earth 的 $5.5$--$10.4$ 倍\,[55]，在被发现时因可能处于其恒星的宜居带而引起关注\,[56]，但后续研究表明它实际上距离其恒星过近，无法宜居。[57]更大质量的行星亦已被发现，其性质与质量向上连续延伸至褐矮星（brown dwarf）领域。[58]

已知一些系外行星与其母恒星之间的距离远小于太阳系中任意一颗行星与 Sun 的距离。太阳系内距离 Sun 最近的 Mercury 距离为 $0.4\ \mathrm{AU}$，其公转周期为 $88$ 天；但超短周期行星（ultra-short-period planets）可以在不到一天内完成公转。Kepler-11 系统中有五颗行星的轨道周期都短于 Mercury，而且它们的质量都远大于 Mercury。  
另有热木星（hot Jupiters），例如 51~Pegasi~b[42]，紧贴恒星运行，可能因蒸发而变成“冥狱行星”（chthonian planets），也就是仅剩下残核的行星。  
此外，还有系外行星比太阳系中任何行星离其恒星更远。Neptune 距 Sun 为 $30$ AU，公转周期为 $165$ 年；但某些系外行星距离其恒星数千 AU，公转周期超过百万年（如 COCONUTS-2b）。[59]
\subsection{属性}
尽管每颗行星都具有各自独特的物理特征，但它们之间仍然存在一些普遍的共同点。其中一些特征——如行星环或天然卫星——目前仅在太阳系行星中被观测到，而另一些则在系外行星中也普遍存在。[60]
\subsection{动力学特征}
\textbf{轨道（Orbit）}
\begin{figure}[ht]
\centering
\includegraphics[width=8cm]{./figures/0d173cd711ffbba6.png}
\caption{行星海王星的轨道与冥王星轨道的比较。请注意冥王星轨道相对于海王星轨道的拉长（偏心率），以及其相对于黄道面的巨大夹角（轨道倾角）。} \label{fig_Planet_7}
\end{figure}
在太阳系中，所有行星绕日公转的方向都与太阳自转方向一致：从太阳北极上方看为逆时针方向。至少已经发现一颗系外行星（WASP-17b）其公转方向与其恒星自转方向相反。[61] 一个行星绕恒星运行一周的周期称为其恒星周期或“年”。[62] 行星的“年长”取决于它与恒星的距离：距离越远，轨道越长，同时由于恒星引力作用较弱，其运行速度也越慢。

没有任何行星的轨道是完全圆形的，因此其与恒星的距离会在一年内发生变化。行星最接近恒星的位置称为“近日点”（在太阳系中称“近日点”）；最远的位置称为“远日点”。当行星接近近日点时，其速度会增大，因为它将引力势能转化为动能，就像地球上自由落体物体下落时会加速一样。当地行星接近远日点时，其速度会减小，就像向上抛出的物体在接近最高点时速度会逐渐减小一样。[63]

每个行星的轨道由一组轨道要素来描述：
\begin{itemize}
    \item 偏心率（eccentricity）描述行星椭圆（椭圆形）轨道的扁长程度。偏心率低的行星轨道更接近圆形，而偏心率高的行星轨道则更为椭圆。太阳系中的行星及大型卫星的偏心率都相对较低，因此它们的轨道几乎是圆形的。[62] 彗星、许多柯伊伯带天体以及若干系外行星则具有非常高的偏心率，因此其轨道极为椭圆。[64][65]
    \item 半长轴（semi-major axis）给出轨道的尺度。它是椭圆轨道最长直径的一半。此值并不等同于远日点距离，因为没有任何行星的轨道是以恒星为椭圆几何中心的。[62]
    \item 轨道倾角（inclination）表示行星轨道相对于某一参考平面的倾斜程度。在太阳系中，参考平面是地球轨道平面，称为黄道面。对于系外行星，参考平面称为“天空平面”，即垂直于地球观察方向的平面。[66] 太阳系八大行星的轨道都非常接近黄道面；然而，一些较小的天体，如智神星（Pallas）、冥王星（Pluto）和阋神星（Eris），以及彗星的轨道倾角都要大得多。[67] 大型卫星的轨道通常也不太偏离其母行星的赤道，但地球的月球、土星的土卫八（Iapetus）以及海王星的海卫一（Triton）是例外。其中海卫一尤为独特，它是唯一逆行绕行母行星的大型卫星，即其轨道方向与母行星自转方向相反。[68]
    \item 行星轨道穿过参考平面的两个点称为升交点（ascending node）和降交点（descending node）。[62] 升交点经度（longitude of the ascending node）是从参考平面零经度到升交点的角度。近日点参数（argument of periapsis，太阳系中称近日点参数）是从升交点到行星最接近恒星位置之间的角度。[62]
\end{itemize}
\textbf{自转轴倾角}
\begin{figure}[ht]
\centering
\includegraphics[width=8cm]{./figures/9e2bb5a76a5b4400.png}
\caption{} \label{fig_Planet_9}
\end{figure}
